\documentclass[11pt,letterpaper,oneside]{report}
\usepackage[T1]{fontenc}
\usepackage[utf8]{inputenc}
\usepackage{lmodern}
\usepackage[letterpaper,margin=0.92in,headheight=15pt]{geometry}
\usepackage{amsmath,amssymb,amsthm,mathtools}
\usepackage{microtype}
\usepackage[dvipsnames]{xcolor}
\usepackage{graphicx}
\usepackage{booktabs,longtable,array,tabularx}
\usepackage{enumitem}
\usepackage{fancyhdr}
\usepackage{titlesec}
\usepackage[most]{tcolorbox}
\usepackage{listings}
\usepackage{tikz}
\usetikzlibrary{positioning}
\usepackage{hyperref}
\hypersetup{colorlinks=true,linkcolor=MidnightBlue,urlcolor=RoyalBlue,pdfauthor={Independent Mathematics and Engineering Study Suite},pdftitle={Linear Systems and Dynamical Systems Guide}}
\definecolor{ink}{HTML}{162639}
\definecolor{navy}{HTML}{173B61}
\definecolor{blue}{HTML}{2874A6}
\definecolor{teal}{HTML}{007C83}
\definecolor{amber}{HTML}{8B5E16}
\definecolor{pale}{HTML}{EDF4F8}
\definecolor{paleteal}{HTML}{ECF8F5}
\definecolor{paleamber}{HTML}{FFF6E8}
\titleformat{\chapter}[display]{\normalfont\huge\bfseries\color{navy}}{\filleft\fontsize{43}{46}\selectfont\thechapter}{0.4ex}{\titlerule\vspace{1ex}}[\vspace{1ex}\titlerule]
\titleformat{\section}{\normalfont\Large\bfseries\color{navy}}{\thesection}{0.8em}{}
\titleformat{\subsection}{\normalfont\large\bfseries\color{teal}}{\thesubsection}{0.75em}{}
\pagestyle{fancy}\fancyhf{}\fancyhead[L]{\small\color{navy}Differential Equations Engineering Suite}\fancyhead[R]{\small\color{navy}Volume 5 of 8}\fancyfoot[C]{\thepage}\renewcommand{\headrulewidth}{0.3pt}
\setlist[itemize]{leftmargin=*,itemsep=3pt,topsep=5pt}\setlist[enumerate]{leftmargin=*,itemsep=4pt,topsep=5pt}
\setlength{\parskip}{5pt}\setlength{\parindent}{0pt}\setlength{\emergencystretch}{2em}
\newtheorem{theorem}{Theorem}[chapter]
\newtheorem{proposition}[theorem]{Proposition}
\newtheorem{lemma}[theorem]{Lemma}
\theoremstyle{definition}\newtheorem{definition}[theorem]{Definition}
\newtheorem{example}[theorem]{Worked example}
\theoremstyle{remark}\newtheorem{remark}[theorem]{Remark}
\newtcolorbox{bookground}[1][]{enhanced,breakable,colback=pale,colframe=blue,boxrule=0.6pt,title={TEXTBOOK ALIGNMENT},fonttitle=\sffamily\bfseries\small,#1}
\newtcolorbox{modern}[1][]{enhanced,breakable,colback=paleteal,colframe=teal,boxrule=0.6pt,title={MODERN COMPUTATIONAL EXTENSION},fonttitle=\sffamily\bfseries\small,#1}
\newtcolorbox{caution}[1][]{enhanced,breakable,colback=paleamber,colframe=amber,boxrule=0.6pt,title={ASSUMPTION / VALIDATION},fonttitle=\sffamily\bfseries\small,#1}
\newcommand{\dd}{\,\mathrm{d}}
\newcommand{\R}{\mathbb{R}}
\newcommand{\lap}{\mathcal{L}}
\newcommand{\norm}[1]{\left\lVert#1\right\rVert}
\newcommand{\abs}[1]{\left\lvert#1\right\rvert}
\lstdefinestyle{engpy}{language=Python,basicstyle=\ttfamily\footnotesize,keywordstyle=\color{RoyalBlue}\bfseries,commentstyle=\color{gray},stringstyle=\color{teal},backgroundcolor=\color{pale},frame=single,rulecolor=\color{blue!30},breaklines=true,columns=fullflexible,keepspaces=true,showstringspaces=false,tabsize=4,upquote=true}
\lstset{style=engpy}
\begin{document}
\begin{titlepage}
\pagecolor{navy}\color{white}
\vspace*{0.75in}
{\sffamily\bfseries\large MATHEMATICS \& ENGINEERING DOCUMENTATION SUITE\\[12pt]}
{\color{cyan!40}\rule{\textwidth}{2pt}}\\[1.0in]
{\sffamily\bfseries\fontsize{29}{36}\selectfont Linear Systems and Dynamical Systems Guide\par}
\vspace{0.45in}
{\Large Volume 5 of 8}\par\vspace{0.55in}
{\color{cyan!30}\rule{0.50\textwidth}{1pt}}\par\vspace{0.3in}
{\large Independent technical companion aligned with the\par}
{\large 13-chapter organization of Nagle, Saff \& Snider (7th ed.)\par}
\vfill
{\large Original derivations, worked examples, verification and engineering practice\par}\vspace{0.3in}
{\small Prepared October 2026 \quad\textbullet\quad pdf\LaTeX{} compatible\par}
\end{titlepage}
\nopagecolor\color{ink}\pagenumbering{roman}
\chapter*{Reader's guide and source boundaries}\addcontentsline{toc}{chapter}{Reader's guide and source boundaries}
This independent study resource is organized around \emph{Fundamentals of Differential Equations and Boundary Value Problems}, Seventh Edition, by R.~Kent Nagle, Edward~B. Saff, and Arthur~David Snider (Pearson, 2018; ISBN 978-0-321-97710-6). The book's chapter and section numbers are given as \emph{navigation references}, not as claims of reproduced text or proofs.

\begin{bookground}
\textbf{Textbook-grounded scope:} topics, terminology, and chapter/section organization corresponding to the cited textbook. All explanations, derivations, exercise prompts, and diagrams in this volume are newly authored and may emphasize different examples. Check the textbook itself for its exact hypotheses, proofs, exercises, or nomenclature in a particular section.
\end{bookground}
\begin{modern}
\textbf{Distinct extension layer:} Python/NumPy/SciPy workflows, engineering QA and reproducibility practices, automated convergence studies, diagnostic plots, and implementation guidance. These are \emph{not} represented as content supplied by the textbook. API calls are examples and should be verified against the installed package version.
\end{modern}
\begin{caution}
\textbf{Mathematical qualification:} formulas require their stated regularity, domain and boundary/initial conditions. Numerical agreement does not replace a proof; computational outputs are not certified error bounds unless explicitly established.
\end{caution}
\tableofcontents\clearpage\pagenumbering{arabic}
\chapter{Systems of first-order differential equations}
\begin{bookground}Correspondence: Chapter 5 (Introduction to Systems and Phase Plane Analysis), particularly Sections 5.1--5.4, and Chapter 9 (Matrix Methods for Linear Systems), Sections 9.1--9.4.\end{bookground}
\section{State vectors and normal form}
A dynamical system in \emph{normal form} is $\mathbf x'(t)=\mathbf f(t,\mathbf x(t))$. Coupled differential equations describe interacting quantities. For instance, two well-mixed compartments exchanging material satisfy
\[\dot x_1=-k_{12}x_1+k_{21}x_2+q_1(t),\quad \dot x_2=k_{12}x_1-k_{21}x_2+q_2(t).\]
When $q_1=q_2=0$, total inventory obeys $(x_1+x_2)'=0$. This conservation law is both a modeling constraint and a numerical invariant.
\section{Higher-order reduction}
An $n$th-order ODE $y^{(n)}=f(t,y,\ldots,y^{(n-1)})$ becomes an $n$-component first-order system using $x_i=y^{(i-1)}$. This reformulation is standard for matrix analysis and numerical solvers. The initial vector must provide all $n$ independent initial conditions; BVPs distribute constraints across domain endpoints instead.
\section{Equilibria and direction fields}
For autonomous $\mathbf x'=\mathbf f(\mathbf x)$, an equilibrium satisfies $\mathbf f(\mathbf x_*)=0$. A \emph{phase portrait} depicts possible orbits and direction arrows in state space. A trajectory is parameterized by time, while an orbit is the geometric locus; do not confuse multiple speeds along similar orbits with identical solutions.
\chapter{Matrix methods and the fundamental matrix}
\begin{bookground}Correspondence: Chapter 9, Sections 9.2--9.8, including matrices, eigenvalues, complex eigenvalues, nonhomogeneous systems and the matrix exponential.\end{bookground}
\section{Homogeneous constant-coefficient systems}
For $\mathbf x'=A\mathbf x$, the unique IVP solution is $\mathbf x(t)=e^{A(t-t_0)}\mathbf x_0$, where
\[e^{At}=\sum_{k=0}^{\infty}\frac{t^kA^k}{k!}.\]
If $A=V\Lambda V^{-1}$, then $e^{At}=Ve^{\Lambda t}V^{-1}$. This formula requires diagonalizability, but the matrix exponential itself is defined for every square matrix. Jordan blocks give terms proportional to $t^ke^{\lambda t}$.
\begin{example}[Spiral sink]
Let $A=\begin{psmallmatrix}-1&-2\\2&-1\end{psmallmatrix}=-I+2J$ where $J=\begin{psmallmatrix}0&-1\\1&0\end{psmallmatrix}$ and $J^2=-I$. Then
\[e^{At}=e^{-t}\begin{pmatrix}\cos2t&-\sin2t\\\sin2t&\cos2t\end{pmatrix}.\]
All trajectories approach the origin while rotating counterclockwise as time increases.
\end{example}
\section{Nonhomogeneous systems}
For $\mathbf x'=A\mathbf x+\mathbf b(t)$, variation of constants yields
\[\mathbf x(t)=e^{A(t-t_0)}\mathbf x_0+\int_{t_0}^t e^{A(t-s)}\mathbf b(s)\dd s.\]
For variable $A(t)$, use a fundamental matrix $\Phi(t)$ instead of $e^{At}$. In systems arising from physical conservation, check whether forcing components are properly included when computing total mass or energy.
\section{Real versus complex eigenmodes}
A real $2\times2$ matrix with eigenvalues $\alpha\pm i\beta$ has real trajectories with factors $e^{\alpha t}\cos\beta t$ and $e^{\alpha t}\sin\beta t$. Real part $\alpha<0$ gives decaying rotation, $\alpha>0$ increasing rotation, and $\alpha=0$ requires more care (linear neutral centers do not imply nonlinear stability).
\chapter{Phase-plane classification and nonlinear systems}
\begin{bookground}Correspondence: Chapter 5 (\S5.4--5.8) and Chapter 12 (\S12.1--12.3).\end{bookground}
\section{Jacobian linearization}
Let $\mathbf x'=\mathbf f(\mathbf x)$ be continuously differentiable near $\mathbf x_*$. Set $\mathbf u=\mathbf x-\mathbf x_*$. Linearization gives $\mathbf u'=J\mathbf u+O(\norm{\mathbf u}^2)$, with $J=D\mathbf f(\mathbf x_*)$. When no eigenvalue of $J$ lies on the imaginary axis, the equilibrium is \emph{hyperbolic}; its local stability type is captured by the linearized dynamics. At a nonhyperbolic equilibrium, the linearization alone may be inconclusive.
\section{Planar linear equilibria}
For a planar real matrix $A$ define trace $\tau=\operatorname{tr}A$ and determinant $\Delta=\det A$. Characteristic roots satisfy $\lambda^2-\tau\lambda+\Delta=0$. If $\Delta<0$, a saddle results. If $\Delta>0$, $\tau<0$, and no zero-real-part eigenvalue, the origin is asymptotically stable. If $\Delta>0$ and $\tau>0$, it is unstable. Discriminant $\tau^2-4\Delta$ distinguishes real and complex eigenvalues, but defective repeated roots need special analysis.
\section{Competing-species example}
Consider $x'=x(1-x-a y)$, $y'=y(1-y-bx)$ with nonnegative populations. The nullclines are $x=0$ or $1-x-ay=0$, and $y=0$ or $1-y-bx=0$. Interior equilibria satisfy a linear pair with $x=(1-a)/(1-ab)$ and $y=(1-b)/(1-ab)$ when $ab\ne1$, and only biologically admissible if both coordinates are positive. Stability is determined by the Jacobian at each admissible equilibrium; no universal coexistence claim follows without parameter restrictions.
\section{Limit cycles, Poincar\'e maps and chaos}
An isolated periodic orbit is a \emph{limit cycle}. A Poincar\'e return map samples a trajectory on a transversal section and reduces a continuous flow to a discrete iteration. Its derivative at a fixed point helps assess local orbital stability. Complex behavior in a plotted numerical trajectory is not by itself proof of chaos; carefully distinguish transient behavior, numerical artifacts and demonstrated sensitive dependence.
\chapter{Lyapunov stability and energy methods}
\begin{bookground}Correspondence: Chapter 12, Sections 12.4 (Energy Methods), 12.5 (Lyapunov's Direct Method), 12.6 (Limit Cycles), and 12.7 (higher-dimensional stability).\end{bookground}
\section{Stability definitions}
An equilibrium is \emph{Lyapunov stable} if sufficiently close initial data remain close for all future time; \emph{attractive} if nearby solutions converge to it; and \emph{asymptotically stable} if both properties hold. These definitions refer to the nonlinear flow and an appropriate invariant domain, not merely the plotted trajectories.
\section{Direct Lyapunov method}
A differentiable candidate $V(\mathbf x)$ is positive definite near the origin if $V(0)=0$ and $V(\mathbf x)>0$ for $\mathbf x\ne0$. Along a trajectory, $\dot V=\nabla V\cdot\mathbf f$. If $\dot V\le0$, this establishes stability under standard local conditions; if $\dot V<0$ for all nonzero nearby states, asymptotic stability follows under appropriate hypotheses. A negative semidefinite derivative alone generally requires an invariance argument for convergence.
\begin{example}[Damped mass-spring energy]
For $m x''+c x'+kx=0$ with $m,k>0$, $c>0$, take $V(x,v)=\frac12kx^2+\frac12mv^2$. Then $\dot V=-cv^2\le0$. The set $\dot V=0$ has $v=0$; invariance of the dynamics forces $x=0$ for the only trajectory remaining in that set. By LaSalle's invariance principle, the origin is asymptotically stable.
\end{example}
\section{Lyapunov equation for linear systems}
For $\mathbf x'=A\mathbf x$, if $A$ is Hurwitz (all eigenvalues have negative real part), then for each symmetric positive-definite $Q$ there is a unique symmetric positive-definite $P$ solving $A^TP+PA=-Q$. Taking $V=\mathbf x^TP\mathbf x$ gives $\dot V=-\mathbf x^TQ\mathbf x$. This provides both a proof technique and a computational certificate for linearized stability.
\chapter{Applications in the textbook's progression}
\begin{bookground}Chapter 5 applications include interconnected tanks, biomathematics, mass-spring systems, electrical systems and dynamical systems. Chapter 12 also considers competing species and neuron models, including FitzHugh--Nagumo equations.\end{bookground}
\section{Compartmental transport}
Represent transport between $n$ compartments by a matrix with nonnegative off-diagonal transfers and physically valid outflow rates; avoid creating or destroying mass through sign mistakes. Positivity preservation and inventory conservation are domain-specific invariants. In a closed, conservative model, $\mathbf 1^T A=0$ under a column-vector convention.
\section{Coupled oscillators}
Two masses connected by springs satisfy $M\mathbf q''+C\mathbf q'+K\mathbf q=\mathbf F(t)$. For symmetric positive-definite $M$ and $K$, the undamped natural frequencies solve $K\mathbf v=\omega^2 M\mathbf v$. A mass-normalized modal basis makes $\mathbf v_i^TM\mathbf v_j=\delta_{ij}$ and diagonalizes the undamped equations in modal coordinates.
\section{Electrical networks}
Kirchhoff laws combined with capacitor relations $i=C\dot v$ and inductor relations $v=L\dot i$ lead to systems of ODEs or differential-algebraic equations depending on topology. Store energy $E=\sum\frac12 C v^2+\sum\frac12 L i^2$ as a sign/units check where applicable. The circuit's forcing, parasitics and grounding assumptions affect the mathematical model.
\chapter{Computational extension: state-space stability audit}
\begin{modern}This NumPy/SciPy workflow is a modern engineering extension to matrix methods and stability analysis. Numerical eigenvalues are approximations and must be interpreted using conditioning and uncertainty.\end{modern}
\begin{lstlisting}
import numpy as np
from scipy.linalg import expm, solve_continuous_lyapunov
from scipy.integrate import solve_ivp

A = np.array([[-1.0, -2.0], [2.0, -1.0]])
eigvals = np.linalg.eigvals(A)
assert np.max(eigvals.real) < 0
x0 = np.array([1.0, 0.0])
sol = solve_ivp(lambda t, x: A @ x, (0, 5), x0,
                rtol=1e-10, atol=1e-12)
assert sol.success
assert np.linalg.norm(sol.y[:, -1] - expm(5*A) @ x0) < 1e-8

# scipy.linalg solves A @ P + P @ A.T = Q.
# Supply A.T and Q=-I to obtain A.T @ P + P @ A=-I.
P = solve_continuous_lyapunov(A.T, -np.eye(2))
assert np.min(np.linalg.eigvalsh(P)) > 0
assert np.linalg.norm(A.T @ P + P @ A + np.eye(2)) < 1e-10
\end{lstlisting}
\section{Model audit record}
For each nonlinear equilibrium record the Jacobian, eigenvalues, type, physically admissible region, invariant sets, and whether the equilibrium is hyperbolic. When eigendirections are poorly conditioned, report residuals $\norm{A\mathbf v-\lambda\mathbf v}$ and test perturbation sensitivity.
\chapter{Study questions and key insights}
\begin{enumerate}
\item Find the eigenvalues and classify $A=\operatorname{diag}(-4,2)$.
\item For $\mathbf x'=A\mathbf x$ with $A=\begin{psmallmatrix}0&-1\\1&0\end{psmallmatrix}$, verify that $\norm{\mathbf x(t)}$ is constant.
\item State why $\dot V\le0$ does not always imply asymptotic stability.
\item A two-compartment closed transfer system has $x_1'=-kx_1+kx_2$, $x_2'=kx_1-kx_2$. Compute eigenvalues and conserved quantity.
\item Explain what a numerical phase portrait can and cannot establish about chaos.
\end{enumerate}
\section{Answer notes}
The first matrix has a saddle (one negative and one positive eigenvalue). For the rotation matrix, $\frac{d}{dt}\norm{\mathbf x}^2=\mathbf x^T(A+A^T)\mathbf x=0$. Semidefinite $\dot V$ can allow non-decaying invariant orbits. For the two-compartment model, eigenvalues are $0$ and $-2k$ and total amount $x_1+x_2$ is constant. A finite computed orbit may suggest dynamical structure, but numerical errors and transient dynamics can imitate it.

\appendix
\chapter{References and reproducibility}
\begin{enumerate}
\item R.~K. Nagle, E.~B. Saff, and A.~D. Snider, \emph{Fundamentals of Differential Equations and Boundary Value Problems}, 7th ed., Pearson, 2018, ISBN 978-0-321-97710-6. Primary source for the \emph{topic map and terminology}; cited locally by chapter and section.
\item SciPy project, \emph{solve\_ivp} reference, \url{https://docs.scipy.org/doc/scipy/reference/generated/scipy.integrate.solve_ivp.html} (modern extension).
\item SciPy project, \emph{solve\_bvp} reference, \url{https://docs.scipy.org/doc/scipy/reference/generated/scipy.integrate.solve_bvp.html} (modern extension).
\item SciPy project, \emph{eigsh} reference, \url{https://docs.scipy.org/doc/scipy/reference/generated/scipy.sparse.linalg.eigsh.html} (modern extension).
\item Python, NumPy, SciPy, and SymPy project documentation; computational checks require recording package versions and tolerances.
\end{enumerate}
\section*{Verification record template}
\begin{tabularx}{\textwidth}{>{\bfseries}p{0.30\textwidth}X}\toprule
Problem identifier & \rule{0pt}{15pt}\ \\ \midrule
ODE/PDE and domain & \\[8pt]
Initial/boundary conditions & \\[8pt]
Units and parameter values & \\[8pt]
Analytic residual & \\ [8pt]
Grid/step refinement & \\ [8pt]
Independent solver comparison & \\ [8pt]
Version and reproducibility notes & \\ [8pt]\bottomrule
\end{tabularx}
\end{document}
